% unidades/21-condicional-tara.tex
\chapter{🔀 \ruby{たら条件}{たらじょうけん} — Condicional たら}
\unidadheader{21}{たら条件}{〜たら}{El condicional más versátil: secuencia, hipótesis y descubrimiento}

\begin{tcolorbox}[vocabbox, title={\emoji{pushpin} Vocabulario}]
\begin{tabularx}{\linewidth}{llX}
\toprule
\textbf{Japonés} & \textbf{Español} & \textbf{Tipo} \\ \midrule
\vocabitem{当たる}{あたる}{acertar / ganar (lotería)}{v. G1}
\vocabitem{宝くじ}{たからくじ}{lotería}{sustantivo}
\vocabitem{卒業する}{そつぎょうする}{graduarse}{v. suru}
\vocabitem{戻る}{もどる}{volver / regresar}{v. G1}
\vocabitem{気づく}{きづく}{darse cuenta}{v. G1}
\vocabitem{着く}{つく}{llegar}{v. G1}
\vocabitem{開ける}{あける}{abrir}{v. G2}
\bottomrule
\end{tabularx}
\end{tcolorbox}

\subsection*{\emoji{speech-balloon} Frases Clave}
\frase{\ruby{家}{いえ}に\ruby{着}{つ}いたら\ruby{電話}{でんわ}します。}{Cuando llegue a casa, te llamo.}
\frase{もし\ruby{宝}{たから}くじが\ruby{当}{あ}たったら\ruby{旅行}{りょこう}します。}{Si gano la lotería, viajaré.}
\frase{\ruby{開}{あ}けたら\ruby{手紙}{てがみ}があった。}{Al abrirlo, había una carta. (descubrimiento)}
\frase{\ruby{卒業}{そつぎょう}したら\ruby{何}{なに}をしたいですか？}{¿Qué quieres hacer cuando te gradúes?}

\subsection*{\emoji{gear} Gramática}
\patron{Formación: [V/adj/N-pasado] + ら}
  {V-past form + ら (たら/だら)}
  {\ej{\ruby{食}{た}べた → \ruby{食}{た}べたら \quad \ruby{来}{き}た → \ruby{来}{き}たら}{verbos regulares}
   \ej{\ruby{良}{よ}かった → \ruby{良}{よ}かったら \quad \ruby{静}{しず}かだった → \ruby{静}{しず}かだったら}{adjetivos}
   \ej{\ruby{学生}{がくせい}だった → \ruby{学生}{がくせい}だったら}{sustantivos}}

\patron{Usos de たら}
  {(1) secuencia temporal \quad (2) hipótesis \quad (3) descubrimiento}
  {\ej{\ruby{宿題}{しゅくだい}が\ruby{終}{お}わったら\ruby{遊}{あそ}んでいいよ。}{Cuando termines la tarea, puedes jugar. (secuencia)}
   \ej{もし100\ruby{万}{まん}円あったら\ruby{何}{なに}をする？}{Si tuvieras 1 millón de yenes, ¿qué harías? (hipótesis)}
   \ej{\ruby{帰}{かえ}ったら、メッセージがあった。}{Cuando llegué a casa, había un mensaje. (descubrimiento)}}

\comparacion{たら vs と (automático)}
  {たら\\\small secuencial / hipotético\\\small \ruby{家}{いえ}に\ruby{着}{つ}いたら\ruby{電話}{でんわ}する\\\textit{Cuando llegue, llamaré (elección)}}
  {と\\\small resultado automático e inevitable\\\small このボタンを\ruby{押}{お}すと\ruby{電気}{でんき}がつく\\\textit{Al presionar, la luz se enciende (siempre)}}
  {たら permite resultados volitivos (llamaré) y es flexible. と solo se usa para resultados automáticos e inevitables, nunca con verbos volitivos en la cláusula resultado.}

\coloquial{たら más frecuente en habla casual}{
  たら es el condicional más común en conversación diaria. 〜たらどう？ (¿Por qué no...? / ¿Qué tal si...?) es una sugerencia casual muy frecuente: 早く\ruby{寝}{ね}たらどう？}

\culturabox{もし〜たら — Fantasías cotidianas en la conversación}{
  Una forma muy común de iniciar conversación en japonés es el «もし〜たら». ¿Y si ganaras la lotería? ¿Y si pudieras vivir en otro país? Estas preguntas hipotéticas (もしゲームで100万円当たったら？) son un recurso social para hacer pequeña conversación y expresar deseos de forma segura e indirecta.
}

\lectura{もし\ruby{宝}{たから}くじが\ruby{当}{あ}たったら\\[0.4em]
  「もし\ruby{宝}{たから}くじが1\ruby{億}{おく}\ruby{円当}{えんあ}たったら、\ruby{何}{なに}をする？」「まず\ruby{家}{いえ}を\ruby{買}{か}う。そして\ruby{世界}{せかい}\ruby{旅行}{りょこう}をする。」「\ruby{仕事}{しごと}は\ruby{辞}{や}める？」「う〜ん、\ruby{好}{す}きな\ruby{仕事}{しごと}に\ruby{就}{つ}きたいな。\ruby{お金}{おかね}があったら、\ruby{夢}{ゆめ}を\ruby{追}{お}えるから。」}
{Si ganara la lotería\\[0.4em]
  «Si ganarás 100 millones de yenes en la lotería, ¿qué harías?» «Primero compraría una casa. Luego viajaría por el mundo.» «¿Dejarías el trabajo?» «Hmm, me gustaría tener un trabajo que me guste. Si tuviera dinero, podría perseguir mis sueños.»}

\begin{tcolorbox}[ejerciciobox, title={\emoji{pencil} Ejercicios}]
\begin{enumerate}
  \item Forma el condicional たら: \ruby{勉強}{べんきょう}する、\ruby{暑}{あつ}い、\ruby{学生}{がくせい}だ
  \item ¿Qué tipo de uso es? ドアを\ruby{開}{あ}けたら、\ruby{猫}{ねこ}がいた。
  \item Crea un もし〜たら sobre tu sueño personal.
\end{enumerate}
\textbf{Respuestas:}
1) \ruby{勉強}{べんきょう}したら・\ruby{暑}{あつ}かったら・\ruby{学生}{がくせい}だったら \quad
2) Descubrimiento \quad
3) (libre)
\end{tcolorbox}

\begin{tcolorbox}[kanjibox, title={\emoji{mahjong} Kanjis}]
\begin{tabularx}{\linewidth}{cllXX}
\toprule
\textbf{Kanji} & \textbf{On} & \textbf{Kun} & \textbf{Significado} & \textbf{Vocab} \\ \midrule
\kanjirow{当}{とう}{\ruby{あ}{たる}}{acertar}{\ruby{当}{あ}たる / \ruby{当然}{とうぜん}}
\kanjirow{卒}{そつ}{—}{graduarse}{\ruby{卒業}{そつぎょう} / \ruby{卒}{そつ}}
\kanjirow{業}{ぎょう}{—}{oficio}{\ruby{卒業}{そつぎょう} / \ruby{業界}{ぎょうかい}}
\kanjirow{戻}{たい}{\ruby{もど}{る}}{regresar}{\ruby{戻}{もど}る / \ruby{返戻}{へんれい}}
\bottomrule
\end{tabularx}
\end{tcolorbox}
